% LaTeX Canvas: Chapter Outline
\chapter{Automation}
\label{ch:automation}

\section*{Purpose}
Automation is how you turn `I can do this once'' into `we can do this reliably.'' In data science and computing projects, automation reduces errors, makes work reproducible, and enables collaboration by ensuring that the same steps run the same way on every machine. This chapter introduces a practical automation ladder: (1) scripts, (2) task runners and rebuild tools, (3) scheduling, and (4) continuous integration (CI). It also shows how to incorporate AI tools responsibly to speed up routine work without weakening verification.

\section*{Learning objectives}
By the end of this chapter, a reader should be able to:
\begin{enumerate}
\item Identify tasks worth automating and write them as repeatable scripts.
\item Use a task runner / build tool (e.g., \texttt{make}) to create named, repeatable commands.
\item Understand incremental rebuilds (targets and dependencies) and why they save time.
\item Schedule scripts to run automatically (cron on Unix-like systems; Task Scheduler on Windows).
\item Explain what CI is and set up a basic CI workflow that runs on pushes and pull requests.
\item Add common CI checks: formatting, linting, tests, and artifact collection.
\item Manage automation hygiene: logs, exit codes, idempotence, and safe failure.
\item Use AI tools to draft automation artifacts (Makefiles, YAML workflows, docs) while maintaining human verification.
\end{enumerate}

\section*{Running theme: make the correct path the easy path}
If running checks and builds is one command, people will do it. If it is a long checklist, people will skip it.

\section{A mental model: the automation ladder}
\subsection{Level 0: manual commands}
\begin{itemize}
\item You type commands repeatedly.
\item Risk: missed steps, inconsistent runs.
\end{itemize}

\subsection{Level 1: scripts}
\begin{itemize}
\item Put the sequence in a \texttt{.py}, \texttt{.sh}, or \texttt{.ps1} file.
\item Benefit: repeatability and shareability.
\end{itemize}

\subsection{Level 2: task runners and rebuild tools}
\begin{itemize}
\item Named tasks: \texttt{make test}, \texttt{make report}.
\item Rebuild only what changed (dependency graph).
\end{itemize}

\subsection{Level 3: schedulers}
\begin{itemize}
\item Run tasks on a time-based schedule (daily, hourly).
\item Benefit: unattended operation.
\end{itemize}

\subsection{Level 4: CI}
\begin{itemize}
\item Run checks automatically on every push/PR.
\item Benefit: catches integration errors early; creates a shared quality gate.
\end{itemize}

\section{Deciding what to automate (student-friendly heuristics)}
\subsection{Automate if}
\begin{itemize}
\item you repeat the task more than twice,
\item you need the task to be done the same way every time,
\item mistakes are costly (grading, publication, deployment),
\item teammates need a one-command workflow.
\end{itemize}

\subsection{Do not automate (yet) if}
\begin{itemize}
\item the process is still unclear,
\item the task changes every hour,
\item automation would hide important learning.
\end{itemize}

\subsection{The ``minimum viable automation'' set}
\begin{itemize}
\item One command to set up environment.
\item One command to run a smoke test.
\item One command to build core outputs.
\end{itemize}

\section{Automation fundamentals: scripts that behave well}
\subsection{Exit codes and failure}
\begin{itemize}
\item Success should exit with code 0.
\item Failure should exit nonzero and print a clear message.
\end{itemize}

\subsection{Idempotence and safe re-runs}
\begin{itemize}
\item A task is idempotent if re-running it yields the same final state.
\item Avoid scripts that require manual cleanup to rerun.
\end{itemize}

\subsection{Inputs/outputs as first-class}
\begin{itemize}
\item Treat file paths and parameters as explicit.
\item Write outputs to predictable locations.
\item Never overwrite raw data.
\end{itemize}

\subsection{Logging and observability (intro)}
\begin{itemize}
\item Print progress markers for long tasks.
\item Save logs for scheduled/CI runs.
\end{itemize}

\section{Repeatable tasks with \texttt{make} (and the idea of rebuilds)}
\subsection{Why use a task runner/build tool}
\begin{itemize}
\item Memorability: \texttt{make lint} beats ``remember five commands.''
\item Standardization: everyone runs the same targets.
\item Incremental rebuilds: do less work when inputs haven't changed.
\end{itemize}

\subsection{Makefile concepts}
\begin{description}
\item[Target] What you want (e.g., \texttt{report.pdf}).
\item[Prerequisites] Inputs the target depends on.
\item[Recipe] Commands to build/update the target.
\end{description}

\subsection{Phony targets for ``commands''}
\begin{itemize}
\item Targets that do not correspond to files (e.g., \texttt{test}, \texttt{clean}).
\item Encourages a consistent interface.
\end{itemize}

\subsection{A student target set (recommended)}
\begin{itemize}
\item \texttt{make setup} (create env, install deps)
\item \texttt{make format} (auto-format code)
\item \texttt{make lint} (static checks)
\item \texttt{make test} (unit tests / smoke tests)
\item \texttt{make run} (execute pipeline)
\item \texttt{make report} (generate figures/tables)
\item \texttt{make clean} (remove generated artifacts)
\end{itemize}

\subsection{Incremental builds for data pipelines (intro pattern)}
\begin{itemize}
\item Raw data $\rightarrow$ processed data $\rightarrow$ figures $\rightarrow$ report.
\item Each step becomes a target with prerequisites.
\item Benefit: change one input and only downstream targets rebuild.
\end{itemize}

\section{Scheduling scripts}
\subsection{Scheduling on macOS/Linux: cron (concept + safe usage)}
\begin{itemize}
\item Cron runs commands when time fields match.
\item Teach: absolute paths, explicit environment, redirected logs.
\item Common novice bug: script works interactively but fails under cron due to PATH/env.
\end{itemize}

\subsection{Scheduling on Windows: Task Scheduler (concept + CLI)}
\begin{itemize}
\item Tasks can run periodically or at specific times.
\item Teach: ``Start in'' directory, user account, and logging.
\end{itemize}

\subsection{Scheduling hygiene}
\begin{itemize}
\item Write logs to a known folder.
\item Add timestamps to outputs.
\item Avoid running overlapping copies (lock files / ``already running'' check).
\item Alerting concept: what should happen if the job fails.
\end{itemize}

\section{Rebuilds and repeatable workflows beyond \texttt{make}}
\subsection{Task runners as interfaces}
\begin{itemize}
\item The main value is a stable command interface.
\item Even if you do not use full dependency graphs, a task runner reduces friction.
\end{itemize}

\subsection{Artifacts vs caches (concept)}
\begin{itemize}
\item Cache: speed up by reusing dependencies.
\item Artifact: save outputs of a run for inspection/download.
\end{itemize}

\section{Continuous Integration (CI): automation as a quality gate}
\subsection{What CI is (student definition)}
\begin{itemize}
\item Frequent integration into a shared branch.
\item Each integration is verified by an automated build/test.
\item Outcome: detect integration errors early.
\end{itemize}

\subsection{CI events and triggers (GitHub Actions framing)}
\begin{itemize}
\item Run on push and pull requests.
\item Use path filters to avoid unnecessary runs.
\end{itemize}

\subsection{Jobs, steps, runners}
\begin{itemize}
\item A workflow is a set of jobs.
\item Jobs run on runners (virtual machines).
\item Steps are commands or reusable actions.
\end{itemize}

\subsection{A minimal CI workflow for students}
\begin{enumerate}
\item Check out code.
\item Set up language/runtime.
\item Install dependencies.
\item Run format/lint checks.
\item Run tests (or smoke tests).
\item Save logs/artifacts if failures occur.
\end{enumerate}

\subsection{CI performance hygiene}
\begin{itemize}
\item Dependency caching to speed up installs.
\item Matrix builds (optional): test multiple versions.
\item Keep CI fast so it runs often.
\end{itemize}

\subsection{CI security basics (novice-appropriate)}
\begin{itemize}
\item Do not print secrets in logs.
\item Understand the difference between untrusted PRs and trusted branches.
\item Use least privilege for tokens and permissions.
\end{itemize}

\section{Incorporating AI tools into automation (responsibly)}
\subsection{What AI is good for}
\begin{itemize}
\item Drafting boilerplate: Makefiles, CI YAML, doc templates.
\item Suggesting test cases and edge checks.
\item Summarizing PRs and highlighting risk areas.
\end{itemize}

\subsection{What AI is not good for (without verification)}
\begin{itemize}
\item Inventing correct tool flags or syntax without testing.
\item Making security decisions (permissions, secrets handling).
\item ``Fixing'' failing CI by random changes.
\end{itemize}

\subsection{Guardrails for AI-assisted automation}
\begin{itemize}
\item Always run locally before trusting CI changes.
\item Use small diffs: one change per PR.
\item Require ``how to test'' steps in PR descriptions.
\item Treat AI output as a draft; confirm against docs.
\item Never paste secrets into prompts.
\end{itemize}

\subsection{A practical workflow: AI as a junior assistant}
\begin{enumerate}
\item You specify the intent and constraints.
\item AI drafts the file (Makefile/workflow).
\item You validate against documentation.
\item You run smoke tests.
\item You open a PR and request review.
\end{enumerate}

\section{Common failure modes and fixes}
\subsection{``Works on my machine'' automation}
\begin{itemize}
\item Fix: pin dependencies, record environment, run in clean env/CI.
\end{itemize}

\subsection{Scheduling failures (wrong PATH, wrong working directory)}
\begin{itemize}
\item Fix: use absolute paths, set working directory, log environment.
\end{itemize}

\subsection{CI failures due to missing files or secrets}
\begin{itemize}
\item Fix: ensure CI does not depend on local files; store non-secret config in repo.
\end{itemize}

\subsection{Slow CI that everyone ignores}
\begin{itemize}
\item Fix: reduce scope, add caching, split jobs.
\end{itemize}

\subsection{Automation that overwrites important artifacts}
\begin{itemize}
\item Fix: write outputs to dedicated folders; add ``are you sure'' safeguards for destructive steps.
\end{itemize}

\section{Worked examples (outline)}
\subsection{Example 1: Turn a 6-step checklist into \texttt{make} targets}
\begin{itemize}
\item Create \texttt{make format}, \texttt{make test}, \texttt{make report}.
\item Add a default target \texttt{make all}.
\end{itemize}

\subsection{Example 2: Schedule a daily pipeline run}
\begin{itemize}
\item Create a script with logs.
\item Add a cron entry (macOS/Linux) or Task Scheduler entry (Windows).
\item Verify outputs and failure behavior.
\end{itemize}

\subsection{Example 3: Add GitHub Actions CI}
\begin{itemize}
\item Run on push and pull request.
\item Install deps, run tests.
\item Upload artifacts on failure.
\end{itemize}

\subsection{Example 4: Speed up CI with caching}
\begin{itemize}
\item Add dependency caching.
\item Compare runtimes before/after.
\end{itemize}

\subsection{Example 5: Add pre-commit hooks + CI enforcement}
\begin{itemize}
\item Run lightweight checks before commit.
\item Ensure CI runs the same checks.
\end{itemize}

\subsection{Example 6: Use AI to draft a workflow, then validate}
\begin{itemize}
\item Draft YAML with AI.
\item Verify triggers/permissions.
\item Run locally and via PR.
\end{itemize}

\section{Templates}
\subsection{Template A: Makefile skeleton (task interface)}
\begin{verbatim}
.PHONY: help format lint test run report clean

help:
@echo "make format | lint | test | run | report | clean"

format:
...

lint:
...

test:
...

run:
...

report:
...

clean:
...
\end{verbatim}

\subsection{Template B: Cron entry (conceptual)}
\begin{verbatim}

# minute hour day month weekday  command

# redirect stdout/stderr to a log file

\end{verbatim}

\subsection{Template C: Windows scheduled task (conceptual)}
\begin{verbatim}

# schtasks /create ...

# run a script at a scheduled time with a named task

\end{verbatim}

\subsection{Template D: GitHub Actions workflow skeleton}
\begin{verbatim}
name: CI
on:
push:
pull_request:

jobs:
test:
runs-on: ubuntu-latest
steps:
- uses: actions/checkout@v4
- name: Set up runtime
...
- name: Install deps
...
- name: Run checks
...
\end{verbatim}

\subsection{Template E: PR checklist for automation changes}
\begin{verbatim}

* What problem does this automation solve?
* How do I test it locally?
* What does it do on failure?
* Any permissions/secrets involved?
* Links to documentation
  \end{verbatim}

\section{Exercises}
\begin{enumerate}
\item Identify three repeated tasks in your project and turn them into \texttt{make} targets.
\item Create a script that writes logs and exits nonzero on failure.
\item Schedule the script (cron or Task Scheduler) and confirm it runs.
\item Add a CI workflow that runs tests on every pull request.
\item Add dependency caching and measure CI runtime change.
\item Add pre-commit hooks; mirror the same checks in CI.
\item Use an AI tool to draft a workflow file, then validate it against official docs and run a test PR.
\end{enumerate}

\section{One-page checklist}
\begin{itemize}
\item I can turn multi-step tasks into one-command targets.
\item My scripts have clear inputs/outputs and safe re-run behavior.
\item Scheduled jobs use explicit paths, working directory, and logging.
\item CI runs on pushes/PRs and includes a minimal quality gate.
\item CI is fast enough to be used continuously; caching is used appropriately.
\item Automation changes are documented and reviewed like code.
\item AI assistance is used to draft, not to bypass verification.
\end{itemize}

\appendix
\section{Quick reference: common automation concepts}
\begin{itemize}
\item Targets, prerequisites, recipes (rebuild logic).
\item Schedulers: cron fields and Windows task schedules.
\item CI: triggers, runners, jobs, artifacts, caches.
\item Hygiene: idempotence, exit codes, logs, secrets.
\end{itemize}
